\subsection{Matter Sector: Neutron Star Dataset (3D Domain)}
\label{sec:neutron-stars}

For static neutron stars the matter is a perfect fluid, $T^{\mathrm{matter}}_{\mu\nu} = (\rho + p)\, u_\mu u_\nu + p\, g_{\mu\nu}$, in hydrostatic equilibrium~\citep{Tolman1939,OppenheimerVolkoff1939},
\begin{equation}
  \frac{dp}{dr} = -(\rho + p) \frac{d\Phi}{dr},
  \label{eq:hydrostatic}
\end{equation}
in the metric~\eqref{eq:metric}, with the pressure and the energy density tied by a single polytropic equation of state (Section~\ref{sec:eda-ns}).~\citep{TODO_EquationOfState} The coupling function of the neutron-star branch is $f(\phi) = \frac{1}{2\beta} \left[ \exp(-\beta \phi^2) - 1 \right]$~\citep{Doneva2017_NS}, so $\epsilon = -1$ in Eq.~\eqref{eq:coupling-conditions}, opposite to the black-hole coupling, for which $\epsilon = +1$.\review{In the revised structure $\lambda$ is the Gauss--Bonnet coupling constant (Section~\ref{sec:action}); the earlier outline called it an EOS parameter -- confirm which the dataset's $\lambda$ is.}

\paragraph{Centre and infinity.}
Regularity at the stellar centre requires $\Lambda(0) = 0$, $\Phi'(0) = 0$ and $\phi'(0) = 0$~\citep{Doneva2017_NS}; the central value $\Lambda_2 = (d^2\Lambda / dr^2)|_{r=0}$ then satisfies a fourth-order algebraic constraint that depends on the central energy density $\rhoc$ and the central pressure $p_0$~\citep{Doneva2017_NS}. Numerical shooting from $r = 0$ across the stellar surface $r = R$ to spatial infinity yields the gravitational mass $M$ and the dilaton charge $\Dch$ through the asymptotics~\eqref{eq:asymptotics}, $\phi(r) \approx \Dch / r + \order{r^{-2}}$~\citep{Doneva2017_NS}.

\paragraph{The dataset.}
The neutron-star dataset~\citep{Doneva2017_NS} maps the three inputs $(\beta, \lambda, \rhoc)$ to the two observables $(M, \Dch)$. Its raw file holds $\nsEdaRawRows$ configurations in $\nsEdaRawRuns$ runs, one per $(\beta, \lambda)$ curve. Two kinds are removed before any analysis: the $\nsEdaEmptyRuns$ runs whose shooting did not converge, and on every curve the unstable configurations past the maximum mass, where $dM/d\rhoc < 0$.~\citep{TODO_MaxMassStability} That leaves $\nsEdaRows$ configurations on $\nsEdaCurves$ curves, each ending at its $\Mmax$; Section~\ref{sec:eda-ns} analyses them. The dataset lists the branch with positive charge, $\Dch > 0$ (Section~\ref{sec:action}), which its preparation checks.
